% Capitulo: Testes e qualidade de dados
\chapter{Testes e qualidade de dados}
\label{cap:testes}

Todo time de dados já passou por isto: um número errado chega a um painel, alguém da área
percebe antes do time, e a confiança em todos os outros números cai junto. Testes de dados
existem para inverter essa ordem, fazendo o time descobrir o problema antes do usuário. Este
capítulo trata do que testar, de contratos entre quem produz e quem consome dados, e de uma
mudança que considero das mais importantes trazidas pela IA ao trabalho com dados: escrever
os testes antes dos modelos deixou de ser uma boa intenção e pode virar requisito.

\section{O que testar}
\label{sec:o-que-testar}

Um teste de dados é uma verificação automática que roda a cada carga e falha quando uma regra
é violada. Alguns tipos cobrem a maior parte dos problemas:

\begin{itemize}
  \item \textbf{Unicidade.} A chave de cada tabela é única. Se o modelo deveria ter uma
    linha por cliente e tem duas para alguns, toda soma por cliente está errada.
  \item \textbf{Não nulos.} Campos obrigatórios estão preenchidos. Uma venda sem data ou sem
    cliente não deveria chegar à camada de consumo.
  \item \textbf{Integridade referencial.} Toda chave que aponta para outra tabela encontra
    correspondência. Uma venda de um produto que não existe na dimensão de produtos some dos
    relatórios por categoria sem avisar ninguém.
  \item \textbf{Valores aceitos.} Campos com lista fechada de valores (status, tipo, região)
    só contêm valores conhecidos.
  \item \textbf{Regras de negócio.} As regras que a área usa para conferir os números: a
    receita líquida nunca é maior que a bruta, o saldo de estoque nunca é negativo, o total
    do modelo bate com o total da fonte.
  \item \textbf{Atualização e volume.} O dado chegou no prazo, e chegou na quantidade
    esperada. Uma carga que trouxe metade das linhas de sempre passa em todos os outros
    testes e está errada.
\end{itemize}

Os quatro primeiros tipos são genéricos e fáceis de aplicar a qualquer tabela; muitas
ferramentas de transformação já os trazem prontos. Os dois últimos exigem conhecer o negócio
e o comportamento dos dados, e são os que mais pegam problemas reais. Colunas financeiras
merecem atenção especial, com testes de faixa e de soma, porque é nelas que um erro custa
mais caro e é notado mais rápido.

\rascunho{
  Relato: um erro que chegou à produção e que um teste simples teria pego (qual era o erro,
  quanto tempo levou para ser percebido, quem percebeu, qual teste passou a existir depois).
}

\section{Testes antes do modelo}
\label{sec:tdd}

O desenvolvimento guiado por testes, ou TDD, inverte a ordem habitual de trabalho. Em vez de
escrever o código e depois testá-lo, primeiro se escreve um teste que descreve o
comportamento esperado, ele falha porque o código ainda não existe, e então se escreve o
código mínimo para fazer o teste passar \parencite{beck2003}. Na engenharia de software, a
prática tem mais de duas décadas. Em dados, sempre foi recomendada e quase nunca praticada. O
motivo era simples: escrever testes custava tempo que o prazo não dava, e ficava para depois,
que nunca chegava.

A IA mudou essa conta. A partir de uma regra de negócio descrita em texto, um assistente
escreve em segundos os testes correspondentes, sugere casos de borda que ninguém tinha
pensado e monta dados de exemplo para exercitá-los. O custo de escrever testes caiu tanto
que o TDD em dados pode deixar de ser um diferencial e passar a ser requisito do time.

O fluxo fica assim. Antes de escrever o modelo, escrevem-se os testes que descrevem as regras
de negócio: \enquote{toda venda cancelada sai da receita}, \enquote{o saldo nunca fica
negativo}, \enquote{o total do frete bate com o dos pedidos}. Esses testes falham, porque o
modelo ainda não existe. Depois se escreve o modelo até que todos passem.

A consequência mais interessante não é técnica. Os testes viram a especificação executável da
regra de negócio. Em vez de validar com o stakeholder cada versão do modelo, o que exige a
presença dele a cada iteração, a validação acontece uma vez, sobre os casos de teste:
\enquote{se o pedido for cancelado depois de faturado, ele sai da receita do mês do
cancelamento ou do mês do faturamento?} Respondida a pergunta, ela vira teste. Daí em diante,
o time itera sobre o modelo sem depender da agenda do stakeholder, e uma mudança de regra vira
uma mudança de teste, discutida e aprovada antes de qualquer código.

\begin{destaque}
Com a IA, escrever testes ficou barato. Os testes viram a especificação da regra de negócio,
validada uma vez com o stakeholder, e o time passa a iterar sem depender da presença dele.
\end{destaque}

Há dois papéis nessa divisão de trabalho. O da IA é gerar os testes a partir da regra
descrita, sugerir casos de borda e montar dados de exemplo. O do profissional é conferir se
os testes dizem mesmo o que o negócio quer, o que exige conhecer o negócio
(capítulo~\ref{cap:negocio}). E há um cuidado decisivo: testes gerados depois do modelo, a
partir do próprio modelo, tendem a confirmar o que ele já faz, inclusive os erros. É por isso
que a ordem importa. O teste escrito antes, a partir da regra, é uma verificação; o teste
escrito depois, a partir do código, é uma fotografia.

\rascunho{
  Relato: um projeto em que os testes foram escritos antes do modelo, com a ajuda da IA, e o
  efeito no tempo de validação com o stakeholder.
}

\section{Contratos de dados}
\label{sec:contratos}

Muitos problemas de dados não nascem no time de dados, e sim na fonte. Um sistema muda o
formato de um campo, renomeia uma coluna ou passa a registrar algo de outra forma, e o
pipeline quebra, ou pior, continua rodando com dados errados. Contratos de dados são uma
tentativa de prevenir isso formalizando o que quem produz o dado se compromete a entregar a
quem o consome \parencite{jones2023}.

Um contrato costuma definir três coisas. A estrutura: quais campos existem, de que tipo, quais
são obrigatórios. O significado: o que cada campo representa e quais valores são válidos. E o
nível de serviço: com que frequência o dado é entregue e com que prazo. Mudanças que quebram o
contrato precisam ser combinadas antes, com tempo para quem consome se adaptar.

Na prática, um contrato só funciona se for verificado automaticamente. Um documento que
descreve a estrutura esperada, sem um teste que confira a cada carga se ela foi respeitada,
envelhece como qualquer documentação. Por isso os contratos e os testes andam juntos: o
contrato é o acordo entre as partes, e os testes são a forma de saber se ele está sendo
cumprido.

\section{Um teste em SQL}
\label{sec:teste-sql}

A forma mais simples e portável de escrever um teste de dados é uma consulta que devolve as
linhas que violam a regra. Se a consulta não devolve nada, o teste passou. O primeiro exemplo
verifica a unicidade da chave de um modelo de receita por cliente:

\begin{sql}
-- Teste: cada cliente aparece uma única vez no modelo
SELECT id_cliente, COUNT(*) AS ocorrencias
FROM receita_por_cliente
GROUP BY id_cliente
HAVING COUNT(*) > 1;
\end{sql}

O segundo verifica uma regra de negócio que pegaria o erro do frete multiplicado
(\veja{sec:consulta-errada}): o frete somado no modelo precisa ser igual ao frete registrado
nos pedidos.

\begin{sql}
-- Teste: o frete do modelo bate com o frete dos pedidos
WITH origem AS (
    SELECT SUM(valor_frete) AS frete FROM pedidos
),
modelo AS (
    SELECT SUM(frete) AS frete FROM receita_por_cliente
)
SELECT o.frete AS frete_origem, m.frete AS frete_modelo
FROM origem AS o
CROSS JOIN modelo AS m
WHERE ABS(o.frete - m.frete) > 0.01;
\end{sql}

Escrito antes do modelo, esse teste falha até que o modelo calcule o frete no grão certo. É
exatamente o tipo de regra que um stakeholder da área financeira confirma em um minuto e que
protege o modelo para sempre. As ferramentas de transformação como código permitem declarar
testes assim junto aos modelos e rodá-los a cada alteração, o que leva ao assunto do próximo
capítulo.

\begin{leituras}
  \leitura{beck2003}{O livro que sistematizou o desenvolvimento guiado por testes. Escrito
    para software, mas a lógica de escrever o teste primeiro se aplica aos dados.}
  \leitura{jones2023}{Uma introdução prática aos contratos de dados, com atenção à relação
    entre quem produz e quem consome.}
  \leitura{moses2022}{Para levar os testes a uma visão mais ampla de observabilidade e
    confiabilidade dos dados.}
\end{leituras}
